% Exam Template for UMTYMP and Math Department courses
%
% Using Philip Hirschhorn's exam.cls: http://www-math.mit.edu/~psh/#ExamCls
%
% run pdflatex on a finished exam at least three times to do the grading table on front page.
%
%%%%%%%%%%%%%%%%%%%%%%%%%%%%%%%%%%%%%%%%%%%%%%%%%%%%%%%%%%%%%%%%%%%%%%%%%%%%%%%%%%%%%%%%%%%%%%

% These lines can probably stay unchanged, although you can remove the last
% two packages if you're not making pictures with tikz.
\documentclass[11pt]{exam}
\RequirePackage{amssymb, amsfonts, amsmath, latexsym, verbatim, xspace, setspace}
\RequirePackage{tikz, pgflibraryplotmarks}

% By default LaTeX uses large margins.  This doesn't work well on exams; problems
% end up in the "middle" of the page, reducing the amount of space for students
% to work on them.
\usepackage[margin=1in]{geometry}
\usepackage{enumerate}
\usepackage{amsthm}

\theoremstyle{definition}
\newtheorem{soln}{Solution}

% Here's where you edit the Class, Exam, Date, etc.
\newcommand{\class}{Math 407 Section 1}
\newcommand{\term}{Fall 2025}
\newcommand{\examnum}{Exam 3}
\newcommand{\examdate}{April 16, 2025}
\newcommand{\timelimit}{1 Hour 50 Minutes}
\newcommand{\ol}[1]{\overline{#1}}

% For an exam, single spacing is most appropriate
\singlespacing
% \onehalfspacing
% \doublespacing

% For an exam, we generally want to turn off paragraph indentation
\parindent 0ex

\begin{document} 

% These commands set up the running header on the top of the exam pages
\pagestyle{head}
\firstpageheader{}{}{}
\runningheader{\class}{\examnum\ - Page \thepage\ of \numpages}{\examdate}
\runningheadrule

\begin{flushright}
\begin{tabular}{p{2.8in} r l}
\textbf{\class} & \textbf{Name (Print):} & \makebox[2in]{\hrulefill}\\
\textbf{\term} &&\\
\textbf{\examnum} & \textbf{Student ID:}&\makebox[2in]{\hrulefill}\\
\textbf{\examdate} &&\\
\textbf{Time Limit: \timelimit} % & Teaching Assistant & \makebox[2in]{\hrulefill}
\end{tabular}\\
\end{flushright}
\rule[1ex]{\textwidth}{.1pt}


This exam contains \numpages\ pages (including this cover page) and
\numquestions\ problems.  Check to see if any pages are missing.  Enter
all requested information on the top of this page, and put your initials
on the top of every page, in case the pages become separated.\\

You may \textit{not} use your books or notes on this exam.\\

You are required to show your work on each problem on this exam.  The following rules apply:\\

\begin{minipage}[t]{3.7in}
\vspace{0pt}
\begin{itemize}

%\item \textbf{If you use a ``fundamental theorem'' you must indicate this} and explain
%why the theorem may be applied.

\item \textbf{Organize your work}, in a reasonably neat and coherent way, in
the space provided. Work scattered all over the page without a clear ordering will 
receive very little credit.  

\item \textbf{Mysterious or unsupported answers will not receive full
credit}.  A correct answer, unsupported by calculations, explanation,
or algebraic work will receive no credit; an incorrect answer supported
by substantially correct calculations and explanations might still receive
partial credit.  This especially applies to limit calculations.

\item If you need more space, use the back of the pages; clearly indicate when you have done this.

\item If the problem asks for a proof, be sure to carefully justify your work, including any theorems from class.

Note: you may NOT use a theorem or result from class to prove something when it makes the problem entirely trivial.  If you are unsure whether a particular theorem or result is allowed, just ask!

Unless otherwise stated, all rings will be commutative with identity and all ring homomorphisms will send the identity to the identity.

\end{itemize}

Do not write in the table to the right.
\end{minipage}
\hfill
\begin{minipage}[t]{2.3in}
\vspace{0pt}
%\cellwidth{3em}
\gradetablestretch{2}
\vqword{Problem}
\addpoints % required here by exam.cls, even though questions haven't started yet.	
\gradetable[v]%[pages]  % Use [pages] to have grading table by page instead of question

\end{minipage}
\newpage % End of cover page

%%%%%%%%%%%%%%%%%%%%%%%%%%%%%%%%%%%%%%%%%%%%%%%%%%%%%%%%%%%%%%%%%%%%%%%%%%%%%%%%%%%%%
%
% See http://www-math.mit.edu/~psh/#ExamCls for full documentation, but the questions
% below give an idea of how to write questions [with parts] and have the points
% tracked automatically on the cover page.
%
%
%%%%%%%%%%%%%%%%%%%%%%%%%%%%%%%%%%%%%%%%%%%%%%%%%%%%%%%%%%%%%%%%%%%%%%%%%%%%%%%%%%%%%

\begin{questions}

\addpoints

\question[10]\mbox{}
\textbf{TRUE or FALSE!}  Write  TRUE if the statement is true.  Otherwise, write FALSE.  Your response should be in ALL CAPS.  No justification is required.
\begin{enumerate}[(a)]
\item  
Every zero divisor is nilpotent
\vspace{1.3in}
\item  
Any unique factorization domain is a principal ideal domain
\vspace{1.3in}
\item  
Every field is a subfield of the complex numbers
\vspace{1.3in}
\item  
In a principal ideal domain, any two elements have a greatest common divisor
\vspace{1.3in}
\item  
The element $7+\sqrt{5}$ is algebraic over $\mathbb{Q}$
\vspace{1.3in}
\end{enumerate}

\newpage
\question[10]\mbox{}
\begin{enumerate}[(a)]
\item  State the definition of prime ideal of a ring $R$
\vspace{2in}
\item  State the definition of a prime element of $R$
\vspace{2in}
\item  Prove that $a\in R$ is prime if and only if $(a)$ is a prime ideal
\end{enumerate}

\newpage
\question[10]\mbox{}

\begin{enumerate}[(a)]
\item  Write the definition of a maximal ideal.
\vspace{2in}
\item  Suppose that $f: R\rightarrow S$ is a ring homomorphism.  Prove that
$$\ker(f) = \{r\in R: f(r) = 0_S\}$$
is an ideal of the ring $R$.
\vspace{3in}
\item  Prove that if $R$ is a field in (b), then either $f$ is trivial or $f$ is injective.
\end{enumerate}


\newpage
\question[10]\mbox{}

Consider the polynomials ring $\mathbb Z[x]$ of polynomials with integer coefficients, the polynomial $p(x) = x^2+x+1$, and the ideal
$$I = (2,p(x)) = \{ 2f(x) + p(x)g(x): f(x),g(x)\in \mathbb Z[x]\}.$$
\begin{enumerate}[(a)]
\item  
Show that any polynomial $f(x)\in \mathbb Z[x]$ is equivalent to a polynomial of degree at most $1$ modulo $I$.
\vspace{4in}
\item 
Show that any polynomial $f(x)\in \mathbb Z[x]$ is equivalent to  $0$, $1$, $x$, or $x+1$ modulo $I$.
\end{enumerate}

\newpage
\question[10]\mbox{}

Carefully prove that $$\mathbb Q(\sqrt{3},\sqrt{5}) = \mathbb Q(\sqrt{3}+\sqrt{5})$$.

\newpage
\question[10]\mbox{}
Let $K$ be a field extension of $F$ and suppose that $u\in K$.
\begin{enumerate}[(a)]
\item  Define the minimal polynomial of an element $u\in K$ which is algebraic over $F$.
\vspace{2in}
\item  Prove that if $f(x)$ is the minimal polynomial of $u$, then $f(x)$ is irreducible.
\vspace{3in}
\item  Prove that if $u\in K$ and $[K:F]=n$, then the minimal polynomial of $u$ has degree at most $n$.
\end{enumerate}

\newpage
\question[10]\mbox{}
Let $F = \mathbb Q$, $K = \mathbb Q(\sqrt[3]{5})$ and $E = \mathbb Q(\sqrt[3]{5},\sqrt{7})$.
\begin{enumerate}[(a)]
\item  Find a basis for $K$ as an $F$-vector space.  Justify your answer!
\vspace{2in}
\item  Find a basis for $E$ over $K$.  Justify!  You may use without proof the fact that $\sqrt{7}\notin K$.
\vspace{3in}
\item  Find a basis for $E$ over $F$.  Justify!
\end{enumerate}


\end{questions}
\end{document}
